\part{Value, Prices, Accounting, and Debt}
\label{part:IV}

Part IV compares valuations, prices, profit, debt, and accounts as representations, and asks what each preserves of the underlying continuation structure.

\begin{housing}Rent, reserves, and the landlord's reported profit are set beside the maintenance backlog, which no conventional account records as a liability.\end{housing}

\chapter{Valuing Feasible Futures}\label{ch:16}
Develops valuations over trajectories, capability bundles, and contingent opportunities. Monetary valuation is one representation among several. Identifies circumstances in which scalar representations exist and circumstances in which relevant comparisons remain partial.
\chspec{\Pref{P13}.}
\draftnote{Candidate import: the identity $D_{\mathrm{KL}}(P(\cdot\mid B)\,\|\,P)=-\log P(B)$ for $0<P(B)\le 1$ as a measure of probabilistic restriction. It does not by itself give a monetary cost, a physical dissipation, or a conservation law.}
\scaffold

\chapter{Prices as Compressed Coordination Signals}\label{ch:17}
Studies prices as economical representations that permit decentralized coordination while omitting much of the underlying causal structure. Asks which omissions are tolerable for a given decision and which require supplementary records, standards, or institutions. Scarcity prices appear as dual variables, and an exchange-rate graph gives the first obstruction example.
\chspec{\Pref{P22}, \Pref{P24}.}
\draftnote{Candidate import: allocation of finite maintenance capacity, $\max_{r_i\ge 0}\sum_i V_i(r_i)$ subject to $\sum_i r_i\le B$, with the budget an explicit assumption and the shadow price as dual variable. No universal conservation of ambiguity is assumed.}
\scaffold

\chapter{Profit and the Location of Costs}\label{ch:18}
Decomposes financial profitability into revenues, recognized costs, transfers of burden, and changes in maintained capacity. A worked enterprise model distinguishes profit supported by reproducible productivity from profit supported by deferred replacement or unpriced damage.
\chspec{\Pref{P12}.}
\scaffold

\chapter{Monetary Debt and Restoration Obligations}\label{ch:19}
Distinguishes financial liabilities from duties to repair ecological, infrastructural, or social damage. Explores when obligations can be transferred or monetized and when substitution fails because the relevant beneficiary, place, or lost capability matters.
\chspec{\Pref{P29}.}
\scaffold

\chapter{Accounting for Continuation}\label{ch:20}
Proposes supplementary accounts for maintenance backlogs, recovery reserves, critical dependencies, and outstanding disposition duties. These accompany conventional accounts without presuming that every dimension can be aggregated into one number. The practical aim is auditable visibility.
\chspec{\Pref{P11}, \Pref{P19}.}
\draftnote{Candidate case (documentation-level only): a two-ledger design that keeps official spending commitments separate from community observations and reports their divergence instead of averaging them. Distinguish an obligation, a reported action, evidence of performance, and discharge of the obligation as four different institutional events.}
\scaffold
